\documentclass[10pt,a4paper]{amsart}
\usepackage{amsmath,amssymb,amsthm,graphicx}
% Retained preamble for the Stacks Project chapter divisors.tex (Divisors).
% The chapter file pinned-original-source.tex carries only \input{preamble} on its first line, so no
% preamble text lives inside the chapter itself.  The Stacks Project uses one single shared preamble for
% every chapter; the snapshot stages it verbatim in sources/stacks-chow/preamble.tex.  The block below is
% copied mechanically, line for line, from that file's lines 14-18 (the Xy-pic setup) and lines 154-217
% (the native theorem declarations together with the mathematical macros).  Nothing was added, removed or
% redefined.  The wrapper class and the amsmath/amssymb/amsthm/graphicx base packages are supplied by the
% assembler, so the class/\externaldocument part of the shared preamble is not repeated here.
% For commutative diagrams we use Xy-pic
\usepackage[all]{xy}

% We use 2cell for 2-commutative diagrams.
\xyoption{2cell}
% Theorem environments.
%
\theoremstyle{plain}
\newtheorem{theorem}[subsection]{Theorem}
\newtheorem{proposition}[subsection]{Proposition}
\newtheorem{lemma}[subsection]{Lemma}

\theoremstyle{definition}
\newtheorem{definition}[subsection]{Definition}
\newtheorem{example}[subsection]{Example}
\newtheorem{exercise}[subsection]{Exercise}
\newtheorem{situation}[subsection]{Situation}

\theoremstyle{remark}
\newtheorem{remark}[subsection]{Remark}
\newtheorem{remarks}[subsection]{Remarks}

\numberwithin{equation}{subsection}

% Macros
%
\def\lim{\mathop{\mathrm{lim}}\nolimits}
\def\colim{\mathop{\mathrm{colim}}\nolimits}
\def\Spec{\mathop{\mathrm{Spec}}}
\def\Hom{\mathop{\mathrm{Hom}}\nolimits}
\def\Ext{\mathop{\mathrm{Ext}}\nolimits}
\def\SheafHom{\mathop{\mathcal{H}\!\mathit{om}}\nolimits}
\def\SheafExt{\mathop{\mathcal{E}\!\mathit{xt}}\nolimits}
\def\Sch{\mathit{Sch}}
\def\Mor{\mathop{\mathrm{Mor}}\nolimits}
\def\Ob{\mathop{\mathrm{Ob}}\nolimits}
\def\Sh{\mathop{\mathit{Sh}}\nolimits}
\def\NL{\mathop{N\!L}\nolimits}
\def\CH{\mathop{\mathrm{CH}}\nolimits}
\def\proetale{{pro\text{-}\acute{e}tale}}
\def\etale{{\acute{e}tale}}
\def\QCoh{\mathit{QCoh}}
\def\Ker{\mathop{\mathrm{Ker}}}
\def\Im{\mathop{\mathrm{Im}}}
\def\Coker{\mathop{\mathrm{Coker}}}
\def\Coim{\mathop{\mathrm{Coim}}}

% Boxtimes
%
\DeclareMathSymbol{\boxtimes}{\mathbin}{AMSa}{"02}

%
% Macros for moduli stacks/spaces
%
\def\QCohstack{\mathcal{QC}\!\mathit{oh}}
\def\Cohstack{\mathcal{C}\!\mathit{oh}}
\def\Spacesstack{\mathcal{S}\!\mathit{paces}}
\def\Quotfunctor{\mathrm{Quot}}
\def\Hilbfunctor{\mathrm{Hilb}}
\def\Curvesstack{\mathcal{C}\!\mathit{urves}}
\def\Polarizedstack{\mathcal{P}\!\mathit{olarized}}
\def\Complexesstack{\mathcal{C}\!\mathit{omplexes}}
% \Pic is the operator that assigns to X its picard group, usage \Pic(X)
% \Picardstack_{X/B} denotes the Picard stack of X over B
% \Picardfunctor_{X/B} denotes the Picard functor of X over B
\def\Pic{\mathop{\mathrm{Pic}}\nolimits}
\def\Picardstack{\mathcal{P}\!\mathit{ic}}
\def\Picardfunctor{\mathrm{Pic}}
\def\Deformationcategory{\mathcal{D}\!\mathit{ef}}

% --- appended 2026-10-12 (Chief): the project own definitions of the reference/slogan/history
% display environments, copied mechanically from the shared Stacks preamble lines 8-12
% (\usepackage{verbatim} plus the three \newenvironment lines, which the chapter preamble keeps
% verbatim from the project preamble). Nothing added, removed or redefined: these environments
% carry no mathematical content and the project itself typesets them as comments. Without these
% three lines a counted statement whose own text uses them fails with Environment slogan undefined.
\usepackage{verbatim}
\newenvironment{reference}{\comment}{\endcomment}
\newenvironment{slogan}{\comment}{\endcomment}
\newenvironment{history}{\comment}{\endcomment}
\usepackage{cleveref}
\crefname{originaldefinition}{original definition}{original definitions}
\Crefname{originaldefinition}{Original definition}{Original definitions}
\makeatletter
\newlabel{algebra-definition-blow-up}{{052Q}{source}{}{sourcepointer-1}{}}
\newlabel{algebra-definition-blow-up@cref}{{[originaldefinition][0][]052Q}{[1][0][]source}{}{}{}}
\makeatother
\crefname{originalsection}{original section}{original sections}
\Crefname{originalsection}{Original section}{Original sections}
\makeatletter
\newlabel{constructions-section-relative-proj-via-glueing}{{01NM}{source}{}{sourcepointer-2}{}}
\newlabel{constructions-section-relative-proj-via-glueing@cref}{{[originalsection][0][]01NM}{[1][0][]source}{}{}{}}
\makeatother
\title{The blowup over an affine open is the Proj of the Rees algebra}
\author{The Stacks Project Authors}
\date{stacks-project:divisors.tex:local-snapshot-20261012; GNU FDL 1.2 or later}
\begin{document}
\maketitle
\setcounter{subsection}{1}
\noindent\textit{Editorial note.} All delivered mathematical text is the author's own, copied from the pinned Stacks Project chapter file divisors.tex; no statement or proof step was rewritten, reordered or removed.\par\medskip
\noindent\textit{Editorial note.} Carrier window: statement and proof both lie inside the window 8000-10000 granted to this lane by CONTROL 8.83.\par\medskip
\noindent\textit{Editorial note.} The card was recovered after a quality read: the prefilter dropped the candidate for the phrase 'The first statement is clear from the construction of the relative Proj via glueing', which names the construction that gives it, while the affine-cover statement is argued in place.\par\medskip
\noindent\textit{Editorial note.} This lemma is the affine input used by several cards of this pool (for instance L44203, L44207, L44212, L44218), where it enters as a navigation pointer.\par\medskip
\noindent\textit{Editorial note.} Preamble: sources/stacks-divisors/retained-preamble.tex (88 lines, verbatim present).\par\medskip
\noindent\textit{Allowed standard prerequisites (editorial).}\begin{itemize}
\item Construction of the relative Proj by glueing. Constructions, Section constructions-section-relative-proj-via-glueing, used at proof lines 8336-8338. Gives the identification of b\textasciicircum{}\{-1\}(U) with the Proj of the Rees algebra.
\item Definition of the affine blowup algebra. Algebra, Definition algebra-definition-blow-up, used at proof lines 8342-8345. Identifies D\_+(a\textasciicircum{}\{(1)\}) with Spec(A[I/a]).
\end{itemize}\medskip
\ifdefined\hypertarget\hypertarget{sourcepointer-1}{}\fi
\noindent\textit{Source reference.} Original definition 052Q: focused-seeds/AGP/sources/stacks-snapshot-20261011/tags.txt, source lines 6578--6578. Navigation pointer into the Stacks Project; the standard fact it supplies is applied inside the counted proof span and the argument resting on it stays there.\par\medskip
\ifdefined\hypertarget\hypertarget{sourcepointer-2}{}\fi
\noindent\textit{Source reference.} Original section 01NM: focused-seeds/AGP/sources/stacks-snapshot-20261011/tags.txt, source lines 2163--2163. Navigation pointer into the Stacks Project; the standard fact it supplies is applied inside the counted proof span and the argument resting on it stays there.\par\medskip
\begin{lemma}
\label{lemma-blowing-up-affine}
Let $X$ be a scheme. Let $\mathcal{I} \subset \mathcal{O}_X$ be a
quasi-coherent sheaf of ideals. Let $U = \Spec(A)$ be an affine open
subscheme of $X$ and let $I \subset A$ be the ideal corresponding to
$\mathcal{I}|_U$. If $b : X' \to X$ is the blowup of $X$ in $\mathcal{I}$,
then there is a canonical isomorphism
$$
b^{-1}(U) = \text{Proj}(\bigoplus\nolimits_{d \geq 0} I^d)
$$
of $b^{-1}(U)$ with the homogeneous spectrum of the Rees algebra
of $I$ in $A$. Moreover, $b^{-1}(U)$ has an affine open covering by
spectra of the affine blowup algebras $A[\frac{I}{a}]$.
\end{lemma}

\begin{proof}
The first statement is clear from the construction of the relative Proj via
glueing, see Constructions, Section
\ref{constructions-section-relative-proj-via-glueing}.
For $a \in I$ denote $a^{(1)}$ the element $a$ seen as an element of
degree $1$ in the Rees algebra $\bigoplus_{n \geq 0} I^n$.
Since these elements generate the Rees algebra over $A$ we see that
$\text{Proj}(\bigoplus_{d \geq 0} I^d)$ is covered by the affine opens
$D_{+}(a^{(1)})$. The affine scheme $D_{+}(a^{(1)})$ is the spectrum of
the affine blowup algebra $A' = A[\frac{I}{a}]$, see
Algebra, Definition \ref{algebra-definition-blow-up}.
This finishes the proof.
\end{proof}

\end{document}
